\songtitle{Laoi Ghláidial}

\tag*{2026}\tag{sean-nos}\tag{irish-lyrics}\tag{african-empire}\tag{monukae}\tag{abulon}\tag{gladiel}\tag{throne-of-smoke}
\begin{notes}
Irish-language (Gaeilge) translation and adaptation by Claude of \emph{Throne of Smoke}, carrying Gladiel's full arc within the \emph{Monukae} saga into a sean-nós duet arrangement. Retrieved from rataflechera's Suno page (V6, 16 September 2026).
\end{notes}
\begin{style}
Irish epic ballad, female mezzo-soprano and male tenor duet with rapid-fire delivery, sean-nós ornamentation and layered choir vocals; 88 bodhrán pulse, half-time chorus; uilleann pipes drone, low whistle countermelody, West African djembe, kora arpeggios, talking drum accents, frame drum rolls, brass fanfare stabs; plate reverb, tape saturation, sidechain pumping, tragic resolve
\end{style}
\begin{lyrics}
\songpart{Verse 1}
An chéad uair dá bhfaca sí é,\\
a chulaith uasal stróicthe ón gcath;\\
fear breá, is rinne sí brionglóid\\
an oíche sin faoi, gan rath.

Ach rinne sí dearmad go luath air —\\
príosúnach eile ag a hathair,\\
airgead fuascailte óna mhuintir shaibhir,\\
nó daor le díol ar an margadh.

Ní fear é a d'fhanfadh sa Mhumhain,\\
ní fear é a gheobhadh Gláidial.

\songpart{Pre-Chorus}
An dara huair a chonaic sí é, ba ghráin léi é.\\
Dá bharr bhí a baile ina lasair is ina luaith,\\
dá bharr do maraíodh a máthair is a muintir,\\
dá bharr do gabhadh í féin is a deartháir.

\songpart{Chorus}
Níor thogh Gláidial a saol —\\
iníon an fhoghlaí mhara,\\
bean ag cibé fear a roghnódh a hathair,\\
ach éan saor uaibhreach go dtiocfadh an lá.

Níor thogh Gláidial an saol eile seo —\\
an cime, an dalta, an compánach,\\
an bhean taobh thiar de ríchathaoir i gcéin,\\
an bhean le ríchathaoir deataigh dá cuid féin.

\songpart{Verse 2}
Fuair Gláidial í féin i gcás óir,\\
na doirse ar oscailt, ach na mílte míle ón muir,\\
i dtír loiscthe an deiscirt, is turas fada farraige\\
ó bhaile nach raibh ann níos mó.

Dinnéar gach lá le buíon a mharaigh a muintir.\\
Chonaic sí a hathair i slabhraí, míonna ina dhiaidh sin,\\
caite i bpríosún nach ligtí di cuairt a thabhairt air.\\
Tháinig a deartháir le focail dheireanacha a n-athar —\\
focail mhearaí os comhair a mhic féin,\\
an mac a mharaigh é.

\songpart{Pre-Chorus}
Troscadh chun báis — ba é sin an saol a thogh Gláidial.\\
Ansin tháinig Abulón, a gabhálaí, chun cainte léi:\\
ní mar leanbh, ní mar mhaighdean, ach mar chomhionann.\\
Athair an fhir ba chúis lena cinniúint\\
ag tairiscint teagasc a thabhairt di —\\
ní mar chime, ní mar iníon, ach mar pháirtí.

\songpart{Chorus}
Níor thogh Gláidial a saol —\\
iníon an fhoghlaí mhara,\\
bean ag cibé fear a roghnódh a hathair,\\
ach éan saor uaibhreach go dtiocfadh an lá.

Níor thogh Gláidial an saol eile seo —\\
an cime, an dalta, an compánach,\\
an bhean taobh thiar de ríchathaoir i gcéin,\\
an bhean le ríchathaoir deataigh dá cuid féin.

\songpart{Bridge}
An croí ba le hathair uair amháin,\\
tá sé ag oscailt anois d'Abulón.\\
Tá an príosún óir ag éirí ina ríocht.

An fear breá ba chúis leis an mbuairt ar fad —\\
Mombutú, mac Abulóin, compánach suairc;\\
agus Marcó, an t-amhas Veinéiseach,\\
an té a dhúisigh an bhean sa chailín seo.

\songpart{Chorus — half}

Níor thogh Gláidial an saol eile seo —\\
an cime, an dalta, an compánach,\\
an bhean taobh thiar de ríchathaoir i gcéin,\\
an bhean le ríchathaoir deataigh dá cuid féin.

\songpart{Verse 3}
Chuaigh Mombutú go dtí a shinsir i bPeiriú, is Marcó leis.\\
Tháinig Mombutú ar ais pósta; níor tháinig Marcó ar ais choíche.\\
Níorbh eol do Ghláidial cé acu fear ba mhó a ghoill uirthi,\\
ach le hais Abulóin d'fhulaingeodh sí.

Ansin dúirt Abulón léi go raibh sí ullamh, a hoiliúint slán:\\
comhairleach do Mhombutú, an bheirt acu i gcúirt na bhFranc.\\
Rogha phraiticiúil, ach urraim fhrithpháirteach leis —\\
faoi dheasghnáth Chríostaí, san Eoraip i gcéin,\\
rinneadh fear céile is bean chéile díobh.

\songpart{Pre-Chorus}
Thiar i dtír loiscthe an deiscirt, Númac, an bhean eile ag Mombutú,\\
chuir sí in iúl go soiléir: ba iad a clann-sean oidhrí Mhonúcae,\\
an talamh ba thorthúla sa deisceart, eastát Abulóin.\\
Ghéill Gláidial. Níor tógadh ise don talamh.

\songpart{Chorus}
Níor thogh Gláidial a saol —\\
iníon an fhoghlaí mhara,\\
bean ag cibé fear a roghnódh a hathair,\\
ach éan saor uaibhreach go dtiocfadh an lá.

Níor thogh Gláidial an saol eile seo —\\
an cime, an dalta, an compánach,\\
an bhean taobh thiar de ríchathaoir i gcéin,\\
an bhean le ríchathaoir deataigh dá cuid féin.

\songpart{Outro}
Fuair Abulón bás ina chath deireanach, i bhfarraigí na Róimhe.\\
Chaoin Gláidial é mar nár chaoin sí a hathair riamh.

Chinntigh an tImpire nach mbeadh cumhacht Abulóin i lámha aon fhir amháin:\\
d'fhág Mombutú Monúcae ag Númac is ag a iníon,\\
faoi chúram a dheartháireacha;\\
d'fhág sé an cabhlach faoi Ghláidial,\\
is é féin ag imirt an taidhleora don Impire.

Tugadh an loingeas trádála suas don Impire —\\
is fágadh loingeas rúnda nua ag Gláidial le tógáil as an nua.

Ríchathaoir deataigh, ach fíorchumhacht.\\
Iníon an fhoghlaí mhara —\\
banríon na mara anois.
\end{lyrics}

